\documentclass[11pt,a4paper]{article}

\usepackage{fontspec}
\setmainfont{Palatino}
\usepackage{geometry}
\geometry{margin=2.5cm}

\usepackage{amsmath}
\usepackage{amssymb}
\usepackage{booktabs}
\usepackage{hyperref}
\hypersetup{colorlinks=true, linkcolor=black, citecolor=black, urlcolor=blue}

% Convenience macros
\newcommand{\Os}{\Omega^{*}}
\newcommand{\G}{\mathrm{G}}
\newcommand{\f}{f}

\sloppy

\title{A Two-Stage Perceptual Account of Interval Consonance}
\author{David De Roure}
\date{Working note, June 2026}

\begin{document}
\maketitle

% ─── Abstract ────────────────────────────────────────────────────────────────

\begin{abstract}
Euler's \emph{Gradus Suavitatis} (1739) assigns a dissonance value to a musical
interval $p/q$ by the formula $\G(p/q) = 1 + \Os(p) + \Os(q)$, where
$\Os(n) = \sum_i e_i(p_i - 1)$ sums the weighted prime exponents of~$n$.
We show that the simpler asymmetric formula $f(p/q) = p + \Os(q)$ achieves
a higher rank correlation with human consonance ratings ($\rho = 0.989$
vs.\ $0.979$), and that this formula admits a natural two-stage perceptual
interpretation grounded in harmonic series theory. We also show that Gradus
is exactly equivalent to a weighted harmonic coincidence count with weights
$w(n) = \Os(n)$, and that the simpler metric $\max(p,q)$ --- equivalent to
Galileo's pulse-coincidence model --- achieves the same correlation as
$f(p/q)$.
\end{abstract}

% ─── 1. Background ───────────────────────────────────────────────────────────

\section{Background}

Euler's Gradus Suavitatis, catalogued in OEIS A275314, assigns to each musical
interval a positive integer measuring its arithmetic dissonance. For a ratio
$p/q$ in lowest terms, the characteristic number is $n = p \cdot q$, and:
\[
  \G(p/q) = 1 + \Os(p \cdot q)
\]
where $\Os$ is the completely additive function $\Os(n) = \sum_i e_i(p_i - 1)$.
Since $\Os$ is completely additive, $\G(p/q) = 1 + \Os(p) + \Os(q)$, treating
numerator and denominator symmetrically. The weights $(p_i - 1)$ for prime
$p_i$ encode Euler's aesthetic judgment that large primes represent greater
``difficulty.''

Euler's formula correctly ranks the 13 standard intervals of Western music
against human consonance ratings \cite{krumhansl1990} with Spearman
$\rho = 0.979$, and is exactly tied only on the major third/major sixth pair
and on a three-way tie at Gradus~8.

% ─── 2. Gradus as harmonic coincidence ───────────────────────────────────────

\section{Gradus as Harmonic Coincidence}

For two notes at frequency ratio $p:q$ (coprime), exact harmonic coincidences
occur at harmonics $p, 2p, 3p, \ldots$ of the lower note. For a fixed count
$M$ of coincidences, the weighted sum with weights $w(n) = \Os(n)$ yields:
\[
  \sum_{k=1}^{M} \Os(k \cdot p) \;=\; \Os(p) \cdot M + C(M)
\]
where $C(M) = \sum_{k=1}^{M} \Os(k)$ is a constant independent of~$p$. This
follows from the complete additivity of $\Os$: $\Os(k \cdot p) = \Os(k) +
\Os(p)$. The symmetric score over both notes gives:
\[
  \text{Score} = M \cdot (\G(p/q) - 1) + 2 \, C(M)
\]

\textbf{Gradus is therefore exactly a weighted harmonic coincidence count},
with harmonics weighted by their own prime complexity. This confirms the
physical intuition behind Euler's formula but reveals its self-referential
structure: recovering Gradus from coincidences requires weights that are
themselves Gradus-like.

Unweighted coincidence counting (flat weights) achieves $\rho = 0.791$ across
all coprime ratios with $p, q \leq 16$. The weight function $\log n$ achieves
$\rho = 0.795$ without any prime structure, establishing a ceiling for
physics-based models that do not invoke prime factorisation.

% ─── 3. Galileo and max(p,q) ─────────────────────────────────────────────────

\section{Galileo's Model and \texorpdfstring{$\max(p,q)$}{max(p,q)}}

Galileo (1638) described consonance in terms of pulse coincidences: two strings
at ratio $p:q$ produce coinciding pulses every $p$ vibrations of the lower
string. The ``sweetness'' of the interval grows with the frequency of
coincidences, i.e.\ inversely with~$p$. The metric $\max(p,q) = p$ is
therefore:

\begin{enumerate}
  \item \textbf{Galileo's pulse model} --- vibrations between coincidences;
  \item \textbf{Harmonic series position} --- the upper note is partial $p$ of
        the implied fundamental \cite{rameau1722};
  \item \textbf{Farey index} --- the smallest $n$ such that $q/p$ appears in
        the Farey sequence $F_n$ (Cauchy 1816);
  \item \textbf{First coincidence harmonic} --- the position in the lower
        note's series where the two notes first meet.
\end{enumerate}

These four descriptions, discovered independently across four centuries, are
mathematically identical. $\max(p,q)$ achieves $\rho = 0.989$ against human
ratings, marginally exceeding Gradus, but cannot distinguish intervals with
the same numerator (e.g.\ $5/4$ and $5/3$ both have $\max = 5$).

% ─── 4. An improved formula ──────────────────────────────────────────────────

\section{An Improved Formula}

Euler's formula is symmetric: $\Os(p) + \Os(q)$. Human ratings, however,
distinguish numerator and denominator --- the bass note and the upper note
play different perceptual roles. We find empirically that the asymmetric
formula:
\[
  \boxed{f(p/q) = p + \Os(q)}
\]
achieves $\rho = 0.989$, matching $\max(p,q)$ and exceeding Gradus. The
formula differs from Gradus only in replacing $\Os(p)$ with~$p$ for the
numerator. Optimising prime weights freely with six parameters (separate
weights for primes 2, 3, 5 in numerator and denominator) achieves perfect
rank correlation $\rho = 1.000$, confirming that the structure required to
perfectly predict the human rankings is present in an asymmetric
prime-weighted formula.

\medskip
\noindent The formula $f(p/q) = p + \Os(q)$ decomposes as the sum of two
quantities:

\medskip
\begin{tabular}{lll}
\toprule
Component & Formula & Meaning \\
\midrule
Upper note reach & $p$ & Partial number of upper note in implied harmonic series \\
Bass context depth & $\Os(q)$ & Prime complexity of lower note's position in series \\
\bottomrule
\end{tabular}

% ─── 5. Two-stage perceptual model ───────────────────────────────────────────

\section{A Two-Stage Perceptual Model}

The formula motivates a two-stage model of how the auditory system evaluates
an interval:

\medskip\noindent
\textbf{Stage 1 --- Bass establishes harmonic context (cost $\Os(q)$).}
The lower note is the $q$-th partial of an implied fundamental $F$. The brain
extrapolates downward to $F$ by ``dividing out'' the prime factors of~$q$.
Each factor of 2 (an octave) costs 1 unit; each factor of 3 costs~2; each
factor of 5 costs~4. This is exactly $\Os(q)$. Intervals whose bass note is a
power of 2 (e.g.\ the octave, $q=1$; the major second bass, $q=8=2^3$)
establish the fundamental cheaply via octave equivalence. Intervals with an
odd prime in $q$ (the fourth, $q=3$; the minor third bass, $q=5$) require the
brain to extrapolate through a genuinely new prime, incurring higher cost.

\medskip\noindent
\textbf{Stage 2 --- Upper note reaches into the series (cost $p$).}
Once the fundamental is established, the upper note must be recognised as its
$p$-th partial. Higher partials are quieter (amplitude $\propto 1/p$ in
typical harmonic sounds), more narrowly tuned, more sensitive to mistuning,
and have lower pitch salience \cite{terhardt1979}. The cost of confirming the
upper note as a member of the harmonic series is therefore $p$ itself.

\medskip\noindent
\textbf{Total perceptual effort:} $f(p/q) = p + \Os(q)$.

\medskip\noindent
This account connects:
\begin{itemize}
  \item \textbf{Galileo} (1638): pulse coincidences as physical mechanism;
  \item \textbf{Rameau} (1722): the \emph{corps sonore} (harmonic series) as
        the source of musical meaning;
  \item \textbf{Euler} (1739): prime arithmetic as the language of harmonic
        complexity;
  \item \textbf{Terhardt} (1979): virtual pitch and harmonic template matching
        as auditory mechanisms;
  \item \textbf{Huron} (2001): statistical learning of harmonic patterns as
        the basis of consonance judgments.
\end{itemize}

% ─── 6. The remaining failure ─────────────────────────────────────────────────

\section{The Remaining Failure}

No simple formula distinguishes the major third ($5/4$, human rank~5) from
the major sixth ($5/3$, human rank~7). Both have $p = 5$ and $\Os(q) = 2$
(since $\Os(4) = \Os(2^2) = 2$ and $\Os(3) = 2$). The distinction requires
knowing that $q = 4 = 2^2$ is simpler than $q = 3$ under octave equivalence,
yet the prime-weighting function cannot see this.

Lahdelma \& Eerola \cite{lahdelma2016} decompose consonance ratings into
sensory and familiarity components; on \emph{sensory consonance} alone,
$5/4$ and $5/3$ are nearly indistinguishable across listeners, suggesting
that the tie in $f(p/q)$ may accurately reflect the purely arithmetic
component of the percept. The gap observed in Krumhansl's data is plausibly
attributable to learned exposure --- major thirds occur far more frequently
than major sixths in tonal music --- consistent with the finding of
McDermott et al.\ \cite{mcdermott2016} that consonance preferences in a
culture with minimal Western music exposure are substantially reduced.
Resolving the tie within an arithmetic formula would therefore require
modelling cultural familiarity, which is outside the scope of the present
approach.

% ─── 7. Summary ──────────────────────────────────────────────────────────────

\section{Summary}

\medskip
\begin{tabular}{lll}
\toprule
Formula & $\rho$ vs.\ human & Notes \\
\midrule
Euler's Gradus: $1 + \Os(p) + \Os(q)$ & 0.979 & Symmetric; two ties unresolved \\
Galileo / $\max(p,q)$                  & 0.989 & Simplest; same tie \\
$f(p/q) = p + \Os(q)$                  & \textbf{0.989} & Asymmetric; two-stage interpretation \\
Optimised asymmetric weights            & 1.000 & Six parameters; no closed form \\
Roughness (Plomp--Levelt)              & 0.709--0.863 & Register-dependent; fails on tritone$^*$ \\
\bottomrule
\end{tabular}

\medskip
\noindent$^*$ Helmholtz \cite{helmholtz1863} introduced the beating-based
dissonance curve that Plomp \& Levelt \cite{plomplevelt1965} later parametrised
from listener data; the two models produce equivalent rankings for our
purposes.

\medskip\noindent
The formula $f(p/q) = p + \Os(q)$ is marginally simpler than Gradus, slightly
more accurate, and directly interpretable as a two-stage process of harmonic
extrapolation. It suggests that Euler's prime-weighting is correct for the
\emph{denominator} but that the \emph{numerator} is better described by the
raw partial number --- consistent with the auditory system's known sensitivity
to partial height rather than prime structure in the upper voice.

% ─── References ──────────────────────────────────────────────────────────────

\begin{thebibliography}{99}

\bibitem{euler1739}
L.~Euler,
\emph{Tentamen novae theoriae musicae},
St.\ Petersburg, 1739.

\bibitem{galilei1638}
G.~Galilei,
\emph{Discorsi e Dimostrazioni Matematiche},
Leiden, 1638.

\bibitem{helmholtz1863}
H.~von Helmholtz,
\emph{On the Sensations of Tone as a Physiological Basis for the Theory of Music},
Vieweg, Braunschweig, 1863.
(English translation by A.~J.~Ellis, Longmans, Green \& Co., London, 1875.)

\bibitem{rameau1722}
J.-P.~Rameau,
\emph{Trait\'{e} de l'Harmonie},
Paris, 1722.

\bibitem{plomplevelt1965}
R.~Plomp and W.~J.~M.~Levelt,
Tonal consonance and critical bandwidth,
\emph{Journal of the Acoustical Society of America} \textbf{38} (1965), 548--560.

\bibitem{terhardt1979}
E.~Terhardt,
Calculating virtual pitch,
\emph{Hearing Research} \textbf{1} (1979), 155--182.

\bibitem{krumhansl1990}
C.~L.~Krumhansl,
\emph{Cognitive Foundations of Musical Pitch},
Oxford University Press, 1990.

\bibitem{lahdelma2016}
I.~Lahdelma and T.~Eerola,
Single chords convey distinct emotional qualities to both na\"{i}ve and
experienced listeners,
\emph{Psychology of Music} \textbf{44} (2016), 37--54.

\bibitem{mcdermott2016}
J.~H.~McDermott, A.~F.~Schultz, E.~A.~Undurraga and R.~A.~Godoy,
Indifference to dissonance in native Amazonian,
\emph{Nature} \textbf{535} (2016), 547--550.

\bibitem{huron2001}
D.~Huron,
Tone and voice: a derivation of the rules of voice-leading from perceptual
principles,
\emph{Music Perception} \textbf{19} (2001), 1--64.

\bibitem{bowling2018}
D.~L.~Bowling, D.~Purves and K.~Z.~Gill,
Vocal similarity predicts the relative attraction of musical chords,
\emph{Proceedings of the National Academy of Sciences} \textbf{115} (2018),
216--221.

\end{thebibliography}

\end{document}
